(1) 上下层自由热应变分别为 \(\alpha_2\Delta T\) 与 \(\alpha_1\Delta T\)，因此 \(\sigma_1(z)=E[\epsilon_0+\kappa z-\alpha_1\Delta T]\)（\(-h<z<0\)），\(\sigma_2(z)=E[\epsilon_0+\kappa z-\alpha_2\Delta T]\)（\(0<z<h\)）。截面无外轴力与外力矩。
(2) \(\int_{-h}^h\sigma b\,dz=0\) 得 \(\epsilon_0=(\alpha_1+\alpha_2)\Delta T/2\)；\(\int_{-h}^h z\sigma b\,dz=0\) 得 \(\kappa=3(\alpha_2-\alpha_1)\Delta T/(4h)\)。正曲率对应上层热伸长更大，且两方程分别消去轴向合力与弯矩。
(3) 令 \(\delta=(\alpha_2-\alpha_1)\Delta T\)，则 \(\sigma_1=E(\delta/2+\kappa z)\)，\(\sigma_2=E(-\delta/2+\kappa z)\)。每单位梁长的弹性储能为 \(U'=b\int_{-h}^h\sigma(z)^2/(2E)\,dz=Eb h\delta^2/16\)，该值非负且与梁长无关。
(4) 若 \(\alpha_1=\alpha_2\)，\(\kappa=0\)、\(\sigma=0\)、\(U'=0\)，恢复无约束自由膨胀。对任意 \(\epsilon_0,\kappa\)，\(U'\) 是正定二次型，零合力零弯矩恰是其驻点条件，故所求弯曲为能量极小而非极大。
